
\begin{table}[H]
\centering
\begin{tabular}{|p{7cm}|p{7cm}|}
\hline
\textbf{Problema} & \textbf{Solução} \\ \hline

Nível de Detalhamento: Alguns sites (ex.: Journey Latin America, Newmarket Holidays) oferecem informações mais granulares, enquanto outros (ex.: Sol Férias, Panam) apresentam estruturas simples. 
& 
Padronização dos campos: criação de estrutura flexível com campos opcionais para dados adicionais e campos obrigatórios para informações essenciais. \\ \hline

Formato de Preços: Alguns sites exibem preços padronizados; outros misturam moeda e texto descritivo. 
& 
Utilização de campos de texto para armazenar valores mistos e posterior normalização via scripts de limpeza. \\ \hline

Descrições e Atributos: Sites europeus apresentam descrições estruturadas (highlights, includes), enquanto sites latino-americanos usam formato livre. 
& 
Uso combinado de campos de texto curto e longo para suportar diferentes níveis de estruturação. \\ \hline

URLs e Navegação: Variação significativa na estrutura das URLs (deep links, hashes, parâmetros complexos). 
& 
Criação de campo URL padronizado e tratamento posterior com scripts de análise para extração de componentes relevantes. \\ \hline

Idioma e Terminologia: Termos específicos em diferentes idiomas (ex.: \textit{circuit/séjour}, \textit{rundreisen}, \textit{holidays}). 
& 
Preservação da terminologia original nos campos de tipo/categoria, garantindo consistência sem perda semântica. \\ \hline

Controle Temporal: Necessidade de rastrear quando os dados foram coletados. 
& 
Inclusão de um campo de data de extração em todas as tabelas. \\ \hline

Heterogeneidade de Campos: Diferenças entre tabelas quanto à presença de colunas específicas. 
& 
Definição mínima comum de colunas (preço, URL, data de extração) para uniformização. \\ \hline

Textos de Diferentes Tamanhos: Campos variando entre descrições curtas e longas. 
& 
Utilização de campos de texto com tamanhos adequados a cada tipo de conteúdo. \\ \hline

\end{tabular}
\caption{Problemas e soluções observados na estrutura dos sites analisados}
\label{tab:problemas_solucoes}
\end{table}